\chapter{线性方程组：解的存在与结构}\label{ch:systems}
\chaptermap{三种计算法}{通解与解空间}{含参数与同解}{相容性与应用}
\begin{examtip}
\textbf{考研核心：}增广矩阵消元、秩判据、基础解系、特解加齐次通解、参数分类。答题不仅要给数值，还应说明自由参数个数。
\textbf{竞赛视角：}把相容性转为列空间与左零空间，把非齐次解集看成仿射平移，用维数控制可能性。
\end{examtip}
\section{高斯消元、克拉默与逆矩阵法}
\subsection{计算方法与可用条件}
\begin{definition}{系数矩阵、增广矩阵与行阶梯形}
方程组记为$A\vx=\vb$，其中$A\in\F^{m\times n},\vx\in\F^n,\vb\in\F^m$；增广矩阵为$[A\mid\vb]$。$\vb=0$时称齐次方程组，$\vb\ne0$时称非齐次方程组；存在解称相容，无解称不相容。行阶梯形满足：全部零行位于非零行之下；每个非零行的首非零元所在列，严格位于上一非零行首非零元的右侧。这个首非零元称主元，所在列称主元列。简化行阶梯形还要求主元全为1，且每个主元列其余元素全为0。相容系统中，系数部分的主元列对应主元变量，其余未知量为自由变量；若增广列出现主元，则该行是$0=c\ne0$，系统无解。
\end{definition}
\begin{theorem}{等价消元与克拉默法则}
初等行变换保持方程组的解集。若$A$是可逆方阵，则唯一解为$A\inv\vb$，并有
\begin{equation}x_j=\dfrac{\det A_j(\vb)}{\det A},\label{eq:cramer}\end{equation}
其中$A_j(\vb)$仅将第$j$列换为$\vb$。消元不要求$A$可逆，克拉默与普通逆矩阵法要求可逆方阵。
\end{theorem}
\begin{proof}
\label{der:4-1}
三类行操作都有逆操作，所以将系数块与右端同时左乘可逆$E$时，$A\vx=\vb$推出$EA\vx=E\vb$，反向再左乘$E\inv$也成立，解集完全相同。若$A$可逆，左乘$A\inv$得到必要形式$\vx=A\inv\vb$；代回成立，因而存在且唯一。

记$A=(a_1,\ldots,a_n)$，唯一解满足$\vb=\sum_i x_i a_i$。把$\vb$置于第$j$列，按该列展开：
\[\det A_j(\vb)=\sum_i x_i\det(a_1,\ldots,\underbrace{a_i}_{j\text{列}},\ldots,a_n).\]
当$i\ne j$，第$i,j$列相同，该项为零；当$i=j$，矩阵就是$A$。所以$\det A_j(\vb)=x_j\det A$，再除以已知非零的$\det A$得克拉默法则。若$\det A=0$，这一步不能相除，算法须返回消元与秩判据。

\end{proof}
\begin{method}{先消元还是先求逆}
单个右端通常直接消元；同一可逆系数配多个右端可复用分解；理论公式或只求一个未知量可用克拉默。易错：初等列变换改变未知量含义，不能在增广矩阵中当作保持原解的行变换使用。
\end{method}
\worked{三法解同一方程组}{基础例题}{\star\star}{4-1}
{解$x+y=3,\ 2x-y=0$。}
{先利用系数少的关系消元，再核对三种公式。}
{\textbf{方法一（消元）：}两式相加，$3x=3$，所以$x=1$；代第一式$y=3-1=2$，代第二式$2-2=0$核验。
\textbf{方法二（克拉默）：}$A=\begin{pmatrix}1&1\\2&-1\end{pmatrix}$，$D=-1-2=-3\ne0$。换第一列得$D_x=\det\begin{pmatrix}3&1\\0&-1\end{pmatrix}=-3$，换第二列得$D_y=\det\begin{pmatrix}1&3\\2&0\end{pmatrix}=-6$，故$x=\frac{D_x}{D}=1,y=\frac{D_y}{D}=2$。
\textbf{方法三（逆）：}$A\inv=\frac13\begin{pmatrix}1&1\\2&-1\end{pmatrix}$，乘右端$(3,0)\T$得$\frac13(3,6)\T=(1,2)\T$。$D\ne0$保证唯一。}
{二阶时三法都短；更高阶通常消元成本最低。}
{克拉默替换的是整列，不是整行。}
\worked{三阶增广消元}{考研风格}{\star\star}{4-2}
{解$x+2y+3z=6,\ 2x+5y+8z=15,\ x+2z=3$。}
{首行可消掉其余行的$x$。}
{增广矩阵及行变换为
\[\left[\begin{array}{ccc|c}1&2&3&6\\2&5&8&15\\1&0&2&3\end{array}\right]
\xrightarrow{\substack{r_2-2r_1\\r_3-r_1}}
\left[\begin{array}{ccc|c}1&2&3&6\\0&1&2&3\\0&-2&-1&-3\end{array}\right]
\xrightarrow{r_3+2r_2}
\left[\begin{array}{ccc|c}1&2&3&6\\0&1&2&3\\0&0&3&3\end{array}\right].\]
末行给$z=1$；第二行$y+2=3$给$y=1$；首行$x+2+3=6$给$x=1$。三个主元$1,1,3$非零，解唯一。代回原三式，左端依次为$6,15,3$，与右端一致。}
{先消元再回代，比计算九个余子式简洁。}
{右端必须同步更新：$15-12=3$、$3-6=-3$。}
\worked{零主元需要换行}{基础例题}{\star}{4-3}
{解$y=2,\ x+y=3$，说明首个主元为零如何处理。}
{把含$x$的方程换到第一行。}
{\textbf{方法一：}增广矩阵$\left[\begin{array}{cc|c}0&1&2\\1&1&3\end{array}\right]$首位为零，交换两行成为$\left[\begin{array}{cc|c}1&1&3\\0&1&2\end{array}\right]$。从末行$y=2$，回代得$x=1$。
\textbf{方法二：}第一式本来已经孤立$y$，直接把$y=2$代入第二式$x+2=3$，得$x=1$。原系数行列式$-1\ne0$，所以零首元仅提示换行，不表示奇异。}
{消元主元为零不代表矩阵奇异；允许换行。}
{换行会改变行列式符号，但不改变方程组解集。}
\worked{克拉默失效后的正确处理}{考研风格}{\star\star}{4-4}
{解$x+y=1,\ 2x+2y=2$。}
{系数行列式为零，改用行消元。}
{第二行减第一行的2倍后，增广矩阵为$\left[\begin{array}{cc|c}1&1&1\\0&0&0\end{array}\right]$。系数秩与增广秩都为1，小于未知量数2。令自由变量$y=t$，主元方程给$x=1-t$，所以
\[\binom xy=\binom10+t\binom{-1}{1},\qquad t\in\R.\]
任意$t$满足$x+y=1$，第二式只是其2倍；反过来每个解都以自己的$y$为$t$进入此式，因此为全部解。$D=0$而非无解，本题不能使用克拉默商式。}
{奇异时消元继续提供秩与自由度。}
{若第二个右端换成3，消元得到$0=1$，此时才是无解。}
\begin{insight}{方程组是在寻找约束的交集}
二元一次方程对应直线。两条直线相交于一点给唯一解；平行且不重合给无解；重合给一条解直线。消元是在保持同一交集的前提下改写约束。秩比较准确地把这种图形分类推广到任意维数。
\end{insight}
\begin{figure}[H]\centering
\begin{tikzpicture}[scale=1]
\begin{scope}
\draw[laaxis](-.2,0)--(2.7,0);\draw[laaxis](0,-.2)--(0,2.7);
\draw[Teal,thick](0,0)--(2.3,2.3);\draw[Gold,thick](0,2)--(2.3,-.3);
\node[lapt]at(1,1){};\node[lanote]at(1.2,-.8){唯一交点\\$r(A)=r(A,\vb)=2$};
\end{scope}
\begin{scope}[xshift=4.2cm]
\draw[laaxis](-.2,0)--(2.7,0);\draw[laaxis](0,-.2)--(0,2.7);
\draw[Teal,thick](0,0)--(2.3,2.3);\draw[Gold,thick](0,.8)--(1.8,2.6);
\node[lanote]at(1.2,-.8){没有交点\\$r(A)<r(A,\vb)$};
\end{scope}
\begin{scope}[xshift=8.4cm]
\draw[laaxis](-.2,0)--(2.7,0);\draw[laaxis](0,-.2)--(0,2.7);
\draw[Teal,line width=2pt](0,0)--(2.3,2.3);\draw[Gold,dashed,thick](0,0)--(2.3,2.3);
\node[lanote]at(1.2,-.8){整条直线\\$r(A)=r(A,\vb)=1$};
\end{scope}
\end{tikzpicture}
\caption{图中$r$表示秩；增广矩阵多出的主元意味着右端与原约束不相容。}\label{fig:system-lines}
\end{figure}
\begin{clevertrick}{参数退化时先想图形，避免盲目除法}
方程$x+y=1,\ tx+y=2$在$t=1$时左边相同、右边不同，对应两条平行线，无法有共同点，所以无解。$t\ne1$时相减得$(t-1)x=1$，此时才可除以$t-1$，得到
\[x=\dfrac1{t-1},\qquad y=1-x=\dfrac{t-2}{t-1}.\]
代回第二式，$tx+y=\frac{t+(t-2)}{t-1}=2$；系数行列式$1-t\ne0$保证交点唯一。图形先指出退化值，代数再给精确坐标。
若把第二个右端改成1，相减改为$(t-1)x=0$。$t=1$时两式重合，全部解为$(1-s,s)$；$t\ne1$时$x=0,y=1$，唯一交点仍存在。相同的系数退化可因右端不同而分别产生无解或无穷多解。
\end{clevertrick}


\section{基础解系与非齐次通解}
\subsection{核空间与仿射平移}
\begin{theorem}{秩判据与解结构}
$A\vx=\vb$有解当且仅当$\rank A=\rank[A\mid\vb]$。若相容且秩为$r$，则$r=n$时唯一，$r<n$时有$n-r$个自由参数。齐次解空间维数为$n-r$；任意一组基称基础解系。非齐次有一个特解$\vx_p$时，全部解为
\begin{equation}\vx=\vx_p+\sum_{j=1}^{n-r}c_j\bm\eta_j,\quad c_j\in\F.\label{eq:general}\end{equation}
\end{theorem}
\begin{proof}
\label{der:4-2}
对$[A\mid\vb]$仅作行变换，先把系数块化为阶梯形。若某个系数全零的行右端为$c\ne0$，对应$0=c$，无解；该右端位置正是增广列新增的主元，所以增广秩比系数秩大。若没有这种行，全部非零系数行均有主元，任给非主元变量，都可从最后一个主元行向上唯一回代，所以有解，增广列也不新增主元。相容性与两秩相等因而等价。

设系数秩$r$，主元变量$r$个，自由变量$n-r$个。齐次时令第$j$个自由变量为1、其他自由变量为0，回代得到$\eta_j$。其自由分量组成单位阵，所以这些向量无关。任意齐次解的自由分量为$c_1,\ldots,c_{n-r}$，减去$\sum_jc_j\eta_j$后自由分量全零，回代使主元分量也全零，因此它们生成全部齐次解，构成基础解系。

相容时选任一特解$\vx_p$。若$\vx$也是解，则$A(\vx-\vx_p)=\vb-\vb=0$，因此$\vx-\vx_p$是上面基础解系的组合。反之，$A(\vx_p+\sum c_j\eta_j)=\vb$，所以式\eqref{eq:general}恰好给全部解。$r=n$时没有自由变量，只有特解；$r<n$时在实数、复数域有无穷多组参数，得到无穷多解。

\end{proof}
\begin{method}{基础解系的参数单位化}
先写出自由变量；每次令其中一个为1、其余为0，回代得到一条基础解向量。特解可令全部自由变量为0。最后检查每个齐次向量都满足原齐次方程，且其自由分量组成单位矩阵。
\end{method}
\worked{两自由变量的基础解系}{考研风格}{\star\star}{4-5}
{求$x_1+2x_2+x_3=0,\ x_2+x_3+x_4=0$的基础解系。}
{令$x_3=s,x_4=t$。}
{两式主元在$x_1,x_2$，令$x_3=s,x_4=t$。第二式给$x_2=-s-t$，第一式给$x_1=-2x_2-x_3=2s+2t-s=s+2t$，所以
\[x=s(1,-1,1,0)\T+t(2,-1,0,1)\T.\]
分别代入两条方向：第一条给$1-2+1=0,-1+1+0=0$；第二条给$2-2+0=0,-1+0+1=0$。若两条的组合为零，看第3分量得$s=0$，看第4分量得$t=0$，故无关。任意解已经由$s,t$表示，所以它们确为基础解系。}
{写出全部参数式并验证独立，才能称基础解系。}
{“找到了两个解”还不够，需证明无关且个数正确。}
\worked{特解加齐次通解}{考研风格}{\star\star}{4-6}
{将例\ref{ex:4-5}右端改为$1,2$，求通解。}
{保留相同齐次方向，只重新求特解。}
{\textbf{方法一：}自由变量设零，第二式给$x_2=2$，第一式$x_1+4=1$给$x_1=-3$，特解$x_p=(-3,2,0,0)\T$。加上例\ref{ex:4-5}的全部齐次方向，得
\[x=(-3,2,0,0)\T+s(1,-1,1,0)\T+t(2,-1,0,1)\T.\]
\textbf{方法二：}直接令$x_3=s,x_4=t$，回代$x_2=2-s-t$、$x_1=1-2(2-s-t)-s=-3+s+2t$，拆开就是上式。代回右端分别为1、2；任意非齐次解减$x_p$落入已求齐次空间，所以没有遗漏。}
{两个方法给同一个仿射平面。}
{特解不唯一，不同特解之差应在齐次解空间中。}
\worked{非齐次解之差}{基础例题}{\star}{4-7}
{若$A\vx_1=A\vx_2=\vb$，证明$\vx_1-\vx_2$为齐次解，并判断和是否仍为原方程解。}
{使用线性性。}
{矩阵乘法对加减线性，因此$A(x_1-x_2)=Ax_1-Ax_2=b-b=0$，差一定是齐次解。对和则$A(x_1+x_2)=2b$。要使它仍解原方程，必须$2b=b$，也就是$b=0$；反向$b=0$时当然满足。因此非齐次右端非零时，和不再属于原解集，而差也不能误当作同右端解。}
{线性解空间与非齐次仿射解集的区别在右端。}
{非齐次解集通常不包含零向量，因此不是线性子空间。}
\worked{解的仿射组合}{竞赛风格}{\star\star}{4-8}
{若$A\vx_i=\vb$，证明$\sum_i c_i\vx_i$在$\sum_i c_i=1$时仍是解。}
{先计算右端，不逐分量展开。}
{\textbf{方法一：}将$A$移入有限线性组合，$A(\sum_i c_ix_i)=\sum_i c_iAx_i=(\sum_i c_i)b=b$，使用的正是系数和为1。
\textbf{方法二：}由$c_1=1-\sum_{i>1}c_i$，组合可写成$x_1+\sum_{i>1}c_i(x_i-x_1)$。每个差满足$A(x_i-x_1)=0$，所以这是一个特解加齐次解，仍为原解。若$b\ne0$，要求组合的右端等于$b$还会反推$\sum_i c_i=1$，所以条件也必要；$b=0$时可不限制系数和。}
{仿射组合保持非齐次解；线性组合还需要系数和条件。}
{$\vb=0$时不必限制系数和；$\vb\ne0$时该限制也是必要的。}
\section{参数分类与同解方程组}
\subsection{不在未知零点上相除}
\begin{definition}{同解与等价约束}
两个方程组同解是指在同一个未知量空间中的解集相同。初等行变换给出同解组，但行数可以不同；新增一个由已有行线性组合得到的方程也不改变解集。
\end{definition}
\begin{theorem}{齐次同解的行空间判据}
实矩阵$A,B$列数相同，则$A\vx=0$与$B\vx=0$同解，当且仅当其行空间相同。对非齐次方程组，仅系数矩阵行空间相同不够，还需右端相容地对应。
\end{theorem}
\begin{proof}
\label{der:4-3}
设$A$的行空间为$U$。$A\vx=0$逐行就是“$\vx$与每一行正交”，也等价于与全部行的线性组合正交，所以$\ker A=U^\perp$。对任意$U\subseteq\R^n$，取其单位正交基并扩充为全空间单位正交基；前面的基生成$U$，其余基恰生成$U^\perp$，再取正交补回到$U$，即$(U^\perp)^\perp=U$。

因此，若$A,B$行空间相同，其核相同，齐次组同解；若齐次组同解，即两核相同，则对两边取正交补，得到两行空间相同。非齐次时还需右端配合，例如$x+y=1$与$x+y=2$的系数行空间相同而解直线不同，这说明增加的条件不能删除。

\end{proof}
\begin{method}{参数分支先列清单}
把消元中将要相除的表达式全部记下；非零分支正常回代；逐个零点重新从未除法的矩阵开始。每个分支同时比较系数秩和增广秩，不只查看$\det A$。
\end{method}
\worked{三角参数系统}{考研风格}{\star\star\star}{4-9}
{解$x+y+z=1,\ ty+z=1,\ (t-1)z=0$。}
{主元可能在$t=0,1$消失。}
{\textbf{分支一：}$t\ne0,1$。末式可除以$t-1$，给$z=0$；第二式$ty=1$给$y=\frac{1}{t}$；首式给$x=1-\frac{1}{t}$，没有自由变量。
\textbf{分支二：}$t=0$。须回到原方程，末式$-z=0$给$z=0$，第二式却为$z=1$，矛盾，故无解。
\textbf{分支三：}$t=1$。末式成为$0=0$，第二式$y+z=1$，首式相减给$x=0$。取$z=s$，得全部解$(0,1-s,s)\T$，$s\in\R$。这三类互斥且覆盖全部参数；特殊值从未代入已经非法相除的中间式。}
{同为行列式零点，一个无解，一个有无穷解。}
{不能先写$z=0$再把$t=1$代入已除过$t-1$的式子。}
\worked{一条左零约束}{考研风格}{\star\star}{4-10}
{讨论$x+y=1,\ 2x+2y=t$。}
{第二行与第一行成比例。}
{\textbf{方法一：}第二式减第一式2倍得$0=t-2$。若$t\ne2$，矛盾无解；若$t=2$，仅剩$x+y=1$，取$y=s$得$(1-s,s)$。
\textbf{方法二：}系数第二行恒为第一行2倍，秩恒1；增广矩阵行变换得到$\left[\begin{array}{cc|c}1&1&1\\0&0&t-2\end{array}\right]$。$t=2$时增广秩1且有一个自由变量；否则增广秩2，与系数秩不等，所以无解。}
{一条必要右端关系即可完成分类。}
{不能因为未知量有两个、方程有两个就默认唯一。}
\worked{不同写法的同一齐次约束}{考研风格}{\star\star}{4-11}
{证明$x+y+z=0,\ x-y=0$与$2x+z=0,\ 2y+z=0$同解。}
{比较行组合或直接参数化。}
{\textbf{方法一：}记第一组左端$g_1=x+y+z,g_2=x-y$，第二组左端$h_1=2x+z,h_2=2y+z$。有$h_1=g_1+g_2,h_2=g_1-g_2$，反向$g_1=\frac{h_1+h_2}{2},g_2=\frac{h_1-h_2}{2}$。所以任一组全为零恰当于另一组全为零，互相推出同解。
\textbf{方法二：}第一组由$x=y=s$得$z=-2s$；第二组先得$x=-\frac{z}{2},y=-\frac{z}{2}$，取$s=-\frac{z}{2}$也为$s(1,1,-2)\T$。两组全部参数集一致，而不只是有一个共同解。}
{互相推出或比较通解均可，行数与外观无关。}
{只验证一个共同非零解不能证明同解。}
\worked{由解集构造方程组}{竞赛风格}{\star\star\star}{4-12}
{构造非齐次方程组，使解集恰为$\vx=e_1+s(e_2-e_1)+t(e_3-e_1)$。}
{两个方向都使分量和为零。}
{将给定向量逐分量写为$(1-s-t,s,t)\T$，其分量和恒为1，所以全部给定向量满足$x_1+x_2+x_3=1$。反过来，若向量满足此式，取$s=x_2,t=x_3$，则$x_1=1-s-t$，它就具有题给形式。因此所求可取单方程$x_1+x_2+x_3=1$。若要多条方程，可加$2x_1+2x_2+2x_3=2$，它只是第一式的倍数，不改变解集；不能添加任何独立约束，否则会缩小解集。}
{先求方向空间的正交补，再用特解确定右端。}
{不能把三个给定点各当成齐次解，右端不能误写为零。}
\section{相容性、左零空间与最小范数}
\subsection{从消元走向空间条件}
\begin{theorem}{相容性的左零空间判据}
实方程组$A\vx=\vb$有解，当且仅当对每个$\vy\in\ker A\T$均有$\vy\T\vb=0$，即
\begin{equation}\im A=(\ker A\T)^\perp.\label{eq:compat}\end{equation}
若有解，解集中存在唯一一个属于$\im A\T$的向量，它是欧氏范数最小的解。
\end{theorem}
\begin{proof}
\label{der:4-4}
若$\vb=A\vx$且$A\T\vy=0$，则$\vy\T\vb=\vy\T A\vx=(A\T\vy)\T\vx=0$，所以$\im A\subseteq(\ker A\T)^\perp$。若$A$有$m$行、秩$r$，则$\dim\ker A\T=m-r$，其正交补维数为$r$，与$\im A$维数相同，故上述包含必为等号。因此“与全部左零向量正交”也充分保证$\vb\in\im A$，即有解。

另由逐行正交条件，$\ker A=(\im A\T)^\perp$。选$\im A\T$的单位正交基并扩充，可将任一特解唯一写成$\vx_p=u+z$，其中$u\in\im A\T,z\in\ker A$。于是$Au=A\vx_p-Az=\vb$，$u$仍是解。任何其他解都为$u+w$且$w\in\ker A$，所以
\[\norm{u+w}_2^2=\norm u_2^2+2u\T w+\norm w_2^2=\norm u_2^2+\norm w_2^2\ge\norm u_2^2.\]
等号当且仅当$w=0$，因此$u$是唯一的最小范数解，也是唯一落在行空间中的解。

\end{proof}
\begin{method}{把多余方程变成右端约束}
求$A\T$的核，令每个基向量与$\vb$内积为零；满足后再回代。若要求最小范数，在通解中令解与所有齐次方向正交。近似最小二乘解留到第7章，它不要求原方程相容。
\end{method}
\worked{超定但相容}{考研风格}{\star\star}{4-13}
{讨论$x=1,y=2,x+y=c$的相容性。}
{第三行是前两行的和。}
{\textbf{方法一：}前两式强制$x=1,y=2$，代第三式要求$c=1+2=3$。若$c=3$，此点满足全部约束且唯一；否则无解。
\textbf{方法二：}$A=\begin{pmatrix}1&0\\0&1\\1&1\end{pmatrix}$，解$A\T y=0$得$y_1+y_3=0,y_2+y_3=0$，所以左零空间为$\Span((-1,-1,1)\T)$。相容性要求$(-1,-1,1)(1,2,c)\T=-1-2+c=0$，仍为$c=3$。$A$前两行含单位阵，满列秩，故相容时唯一。}
{方程数多于未知量不一定无解；关键是右端是否满足依赖。}
{唯一解对应满列秩，不必要求矩阵方形。}
\worked{零乘积与自由度}{竞赛风格}{\star\star}{4-14}
{若$A\in\F^{m\times n},B\in\F^{n\times p}$且$AB=0$，证明$\rank A+\rank B\le n$。}
{把$B$的每列看成齐次解。}
{设$B$的任意列为$b_j$。由$AB=0$逐列知$Ab_j=0$，所以每列都在$\ker A$，进而全部列组合$\im B\subseteq\ker A$。取维数，$\rank B\le n-\rank A$，移项得到结论。另一种写法是从$B$选$r_B$个独立列，它们都是独立齐次解；基础解系只有$n-r_A$维，因此$r_B\le n-r_A$。两个论证都使用相乘处的公共维数$n$，而不是$m$或$p$。}
{零乘积给空间包含，比对元素展开有效。}
{包含于$\ker A$，一般不是$\ker A\T$，后者维数也可能不同。}
\worked{所有右端均有解}{竞赛风格}{\star\star\star}{4-15}
{证明$m\times n$矩阵$A$对每个$\vb\in\F^m$都有解，当且仅当存在右逆$R$满足$AR=I_m$。}
{将单位矩阵的各列作为右端。}
{\textbf{方法一：}若任意右端可解，分别以$e_1,\ldots,e_m$作右端，选出$x_j$满足$Ax_j=e_j$。拼$R=(x_1,\ldots,x_m)$，逐列$AR=I_m$。反向若$AR=I_m$，对任意$b$取$x=Rb$，则$Ax=ARb=b$，故全部右端可解。
\textbf{方法二：}任意右端可解意味着$\im A=\F^m$，秩$m$。选$m$个独立列组成可逆$m$阶矩阵$C$，令$E$为从$I_n$中取相同编号列组成的$n\times m$选择矩阵，则$AE=C$。取$R=EC\inv$，就有$AR=CC\inv=I_m$，给出具体构造。}
{右逆对应满射，左逆对应单射，方阵时二者才合一。}
{长方阵有右逆时未必有左逆，例如$A=(1,0)$。}
\worked{最小范数解}{竞赛风格}{\star\star}{4-16}
{在$x+y=1$的所有解中求$x^2+y^2$最小者。}
{特解应与齐次方向$(1,-1)\T$正交。}
{\textbf{方法一：}解直线的方向是$(1,-1)$。最小范数解与方向正交，故$x-y=0$；联立$x+y=1$得$x=y=\frac{1}{2}$，平方范数$\frac{1}{4}+\frac{1}{4}=\frac{1}{2}$。
\textbf{方法二：}消去$y=1-x$，目标为$x^2+(1-x)^2=2x^2-2x+1=2(x-\frac{1}{2})^2+\frac{1}{2}$，唯一最小点同上。
\textbf{方法三：}Cauchy给$(x+y)^2\le(1^2+1^2)(x^2+y^2)$，由约束得$x^2+y^2\ge\frac{1}{2}$；等号要求$(x,y)$平行于$(1,1)$，配合约束仍为$(\frac{1}{2},\frac{1}{2})$，下界确可达到。}
{正交法可推广高维，配方法直接，Cauchy给最短下界证明。}
{最小范数问题的约束仍要满足，不能把零向量作为候选解。}
\begin{insight}{通解是“一个位置加全部方向”}
对$x+y=1$，向量$(1,0)\T$给出一个位置，齐次解方向$(1,-1)\T$决定整条直线的走向。平移位置不会改变方向空间。原点到这条解直线最近的点，恰与齐次方向正交，因此最小范数问题也能直接从图形得到条件。
\end{insight}
\begin{figure}[H]\centering
\begin{tikzpicture}[scale=2.5]
\draw[laaxis](-.55,0)--(1.65,0)node[right]{$x$};\draw[laaxis](0,-.55)--(0,1.65)node[above]{$y$};
\draw[Teal,thick](-.4,1.4)--(1.4,-.4)node[below right]{$x+y=1$};
\draw[dashed,Gold](-.5,.5)--(.65,-.65)node[below]{$x+y=0$};
\draw[lavec,Ink](0,0)--(.5,.5)node[midway,left]{$\vx_{\min}$};
\draw[lavec,Gold](.5,.5)--(.9,.1)node[midway,above right]{$\ker A$方向};
\node[lapt]at(.5,.5){};\node[lanote,anchor=south west]at(.55,.55){$(\frac{1}{2},\frac{1}{2})$};
\draw[Ink!70](.43,.57)--(.5,.64)--(.57,.57);
\node[lapt]at(1,0){};\node[lanote,anchor=north west]at(1.02,-.02){特解$(1,0)$};
\end{tikzpicture}
\caption{非齐次解直线平行于核；最小范数解是原点到解直线的垂足。}\label{fig:affine-minnorm}
\end{figure}
\begin{clevertrick}{不先写所有参数，也能找最小范数解}
对约束$\va\T\vx=b$且$\va\ne0$，先在法向方向寻找解。设$x_0=c\va$，代入$c\va\T\va=b$，因$\va\T\va>0$得
\[x_0=\dfrac{b}{\va\T\va}\va.\]
为什么它最优？任意解$\vx$与$x_0$之差$h$满足$\va\T h=0$，所以$x_0\T h=0$。勾股分解给
\[\norm\vx^2=\norm{x_0+h}^2=\norm{x_0}^2+\norm h^2\ge\norm{x_0}^2,\]
等号仅在$h=0$时成立，故$x_0$就是唯一最小范数解。对$x+2y+2z=3$，$\va=(1,2,2)\T$，$\va\T\va=1+4+4=9$，所以$x_0=(\frac{1}{3},\frac{2}{3},\frac{2}{3})\T$；约束左端为$\frac{1}{3}+\frac{4}{3}+\frac{4}{3}=3$，长度平方$\frac{1}{9}+\frac{4}{9}+\frac{4}{9}=1$。另一法用Cauchy：$|b|\le\norm\va\norm\vx$给长度下界$\frac{|3|}{3}=1$，而$x_0$确能取到。
\textbf{易错提示：}垂足相对于所选内积而言；带权平方和问题的法向条件要相应改变。
\end{clevertrick}


\begin{warning}
齐次方程总有零解；“无解”只发生于非齐次相容性失败。非齐次解的差为齐次解，和通常不是。基础解系个数是未知量数减秩；与方程条数减秩对应的是左零空间。把任意一个特解写成通解会遗漏全部自由方向。
\end{warning}
\section*{章末分层习题}\addcontentsline{toc}{section}{章末分层习题}
\exercise{4-1}{基础}{解$2x+y=5,x-y=1$，用两法核验。}
\exercise{4-2}{基础}{求$x+2y+z=0$的一组基础解系。}
\exercise{4-3}{基础}{写出$x+2y+z=3$的通解。}
\exercise{4-4}{考研风格}{讨论$x+y=1,tx+y=2$的解。}
\exercise{4-5}{考研风格}{四个未知量的相容方程组系数秩为2，其解集有几个自由参数？非齐次解集是否线性空间？}
\exercise{4-6}{考研风格}{若$\eta_1,\eta_2$是$A\vx=\vb\ne0$的解，求$2\eta_1-\eta_2$与$\eta_1+\eta_2$对应的右端。}
\exercise{4-7}{竞赛拓展}{证明满列秩实矩阵有左逆，并写出一个公式。}
\exercise{4-8}{竞赛拓展}{求$\begin{pmatrix}1&0\\0&1\\1&1\end{pmatrix}\vx=(a,b,c)\T$相容的充要条件。}
\exercise{4-9}{竞赛拓展}{求$x+y+z=3$的最小范数解。}
\section*{章末方法速查}\addcontentsline{toc}{section}{章末方法速查}
\begin{quicktable}
增广消元&任意矩阵&多未知量或奇异系数&右端同步\\
克拉默&可逆方阵&只求一个分量&替换整列\\
基础解系&齐次系统&存在自由变量&个数$n-r$\\
特解加核&非齐次相容&要求全部解&特解不是唯一\\
参数分支&主元依赖参数&行列式零点&不得丢零分支\\
左零判据&超定或多余约束&右端相容性&用$A\T$的核\\
最小范数&相容但不唯一&欧氏距离最小&与核正交\\
\end{quicktable}
